\documentclass[pdflatex,sn-mathphys-num]{sn-jnl}

\usepackage{graphicx}
\usepackage{amsmath,amssymb,amsfonts}
\usepackage{booktabs}
\usepackage{xcolor}

\raggedbottom
\newcommand{\publicrepo}{https://github.com/brunomalhano/vlm-prelabeling-pipeline}

\begin{document}

\title[Prompt Formulation in Vision-Language Pre-labeling]{Prompt Formulation Effects in Modular Vision-Language Pre-labeling Pipelines for Instance Segmentation}

\author*[1]{\fnm{Bruno Malhano de Oliveira} \sur{Jordao}}\email{bruno.42640054@ucp.br}
\author[1]{\fnm{Giovane} \sur{Quadrelli}}\email{giovane.quadrelli@ucp.br}

\affil*[1]{\orgname{Universidade Católica de Petrópolis (UCP)}, \orgaddress{\street{Rua Benjamin Constant, 213, Centro}, \city{Petrópolis}, \state{RJ}, \postcode{25620}, \country{Brazil}}}

\abstract{Pre-labeling with Vision-Language Models (VLMs) can reduce annotation cost in segmentation datasets, but prompt formulation is usually chosen without empirical evidence. We evaluate how prompt formulation affects instance-mask quality in a modular Grounding DINO (Swin-B) + SAM 2.1 (Hiera Large) pipeline on 500 COCO val2017 images across 10 classes.

We compare four formulations of the same class names: simple, direct, contextual, and object-prefixed. Results are counter-intuitive: simple prompts perform best in all aggregation views, reaching mIoU 0.527 [95\% CI: 0.515--0.539] in micro-average, 0.506 [0.489--0.522] in macro-average, and 0.513 [0.477--0.547] under strict class+instance balancing. Alternative formulations are consistently inferior. Differences are statistically significant under paired Wilcoxon signed-rank tests ($p < 10^{-54}$), with medium-to-large effect sizes (Cliff's $\delta$ = 0.43--0.52).

Stage-wise decomposition shows that 73.8\% of failures are grounding misses, while pure segmentation failures account for only \textasciitilde2.2\%. When grounding succeeds (box IoU $\geq$ 0.75), SAM 2.1 attains Mask IoU 0.849, indicating that segmentation quality is sufficient for practical pre-labeling once localization is correct.

These findings provide actionable guidance: use single class-name prompts as default, prioritize optimization of the grounding stage, and perform class-level ROI triage before large-scale pre-labeling deployment.}

\keywords{instance segmentation, pre-labeling, vision-language models, Grounding DINO, SAM 2.1, prompt formulation}

\maketitle

\section{Introduction}

The development of computer vision systems depends on high-quality labeled datasets. In segmentation tasks, annotation requires pixel-level delineation for each object region and is significantly more costly than classification or bounding-box labeling [1].

Foundation VLM pipelines have made automatic pre-labeling practical. Grounding DINO localizes arbitrary objects from text prompts [2], and Segment Anything (SAM) produces masks from visual prompts such as boxes [3]. Combined in a modular pipeline, these models can generate candidate masks without retraining.

Recent work demonstrates feasibility at scale. Labeling Copilot reports agentic data curation with multi-model orchestration [5]. Complementarily, IEEE document 10883307 compares manual, interactive AI-assisted, and pre-labeling-assisted instance annotation regimes [6].

However, a central practical question remains underexplored: which text prompt formulation should practitioners use? Most open-vocabulary pipelines adopt simple class names by convention [2,7], with limited empirical justification.

This study evaluates four prompt formulations in a Grounding DINO + SAM 2.1 pipeline over COCO val2017 and isolates where degradation occurs (grounding vs segmentation). We show a robust \emph{less is more} effect: simple prompts consistently outperform more elaborate variants.

\section{Related Work}

\subsection{Open-vocabulary Segmentation and Grounding}

OVSeg [8] and related open-vocabulary approaches improved text-region alignment, but generally assume canonical class-name prompts. Grounding DINO [2] became a de facto text-to-box component in modular pipelines, and UAV-focused reviews highlight deployment constraints under scale and domain shift [7].

From a deployment perspective, these works establish feasibility, but prompt phrasing itself is typically treated as a fixed design choice rather than an experimental variable.

\subsection{Promptable Segmentation}

SAM [3] introduced promptable segmentation at foundation-model scale. SAM 2 improved robustness and video/image support [4]. In modular pipelines, segmentation quality depends strongly on prompt-box quality, making grounding quality a first-order factor.

\subsection{Assisted Annotation and Prompt Engineering}

Interactive annotation acceleration predates foundation models (e.g., Polygon-RNN++) [9]. More recent agentic systems include Labeling Copilot [5], while IEEE 10883307 reports productivity gains for AI-assisted labeling workflows [6].

For prompt engineering, CLIP established template sensitivity [10], with CoOp and ProDA extending learnable prompt representations [11,12]. RegionCLIP and APE indicate prompt structure matters for region-level grounding [13,14].

The specific gap addressed here is prompt formulation in a fixed modular grounding-to-segmentation pre-labeling stack. Existing studies motivate this factor but do not isolate it under controlled thresholds, data splits, and image-level paired inference.

\section{Methodology}

\subsection{Problem Formulation}

For each image $I$ and class $c$, a prompt template $\tau(c)$ is provided to Grounding DINO, producing candidate boxes. SAM 2.1 then converts accepted boxes into masks. Let predicted instances be $\hat{\mathcal{Y}}=\{(\hat{b}_k,\hat{m}_k)\}_{k=1}^{K}$ and ground-truth instances $\mathcal{Y}=\{(b_j,m_j)\}_{j=1}^{J}$ for the same class in the same image.

\subsection{Pipeline Stages and Prompts}

The pipeline has two stages:
\begin{itemize}
\item \textbf{Grounding}: Grounding DINO (Swin-B) [2] maps text prompts to candidate boxes.
\item \textbf{Segmentation}: SAM 2.1 (Hiera Large) [4] maps boxes to instance masks.
\end{itemize}

Prompt types:
\begin{itemize}
\item \textbf{Simple}: \texttt{dog}
\item \textbf{Direct}: \texttt{segment the dog}
\item \textbf{Contextual}: \texttt{the dog in the image}
\item \textbf{Object-prefixed}: \texttt{object: dog}
\end{itemize}

\begin{figure}[t]
\centering
\fbox{\begin{minipage}{0.92\linewidth}
\centering
Prompt template $\rightarrow$ Grounding DINO (text$\rightarrow$boxes) $\rightarrow$ SAM 2.1 (boxes$\rightarrow$masks) \\
$\rightarrow$ Hungarian matching (prediction$\leftrightarrow$ground truth) $\rightarrow$ metrics and inferential analysis
\end{minipage}}
\caption{End-to-end evaluation workflow.}
\label{fig:pipeline-overview}
\end{figure}

\subsection{Hungarian Matching and Error Definitions}

Matching is performed \textbf{per class per image}. A pairwise mask-IoU matrix is built and converted to cost:

$$
C_{kj}=1-\text{IoU}(\hat{m}_k,m_j).
$$

Optimal bipartite assignment is solved using Hungarian matching [15,16]. Matches below IoU 0.10 are discarded. Unmatched predictions are counted as false positives, and unmatched ground-truth instances as misses. This one-to-one assignment follows DETR-style evaluation principles [17].

\subsection{Metrics and Inferential Analysis}

We report box IoU, mask IoU, Dice, boundary F1, detection rate, and utility buckets (good/correctable/bad). Unmatched misses and false positives are assigned zero mask quality for aggregate reporting.

Inference uses bootstrap 95\% confidence intervals and paired Wilcoxon signed-rank tests with Cliff's $\delta$. To reduce dependence inflation, Wilcoxon tests are computed on \textbf{image-level aggregated mask IoU} (paired across prompts by image).

\section{Experimental Setup}

\subsection{Dataset and Sampling Protocol}

Experiments use COCO val2017 [1] with a stratified sample of 500 images (seed=42) and 10 target classes. The sampling routine enforces at least 30 instances per target class and then fills the sample by prioritizing multi-class coverage.

The resulting EN-only evaluation includes 4590 class-instance records per prompt type.

To support both deployment realism and class fairness, we report two complementary views throughout the analysis: (i) the natural class distribution of the sampled dataset, and (ii) a class-balanced view computed by downsampling each class to the minimum available count (54 instances/class in this run).

\begin{table}[t]
\caption{Per-class instance counts in the analyzed sample (EN run).}
\label{tab:sample-distribution}
\centering
\begin{tabular}{lc}
\toprule
Class & Instances \\
\midrule
person & 1863 \\
chair & 722 \\
bottle & 535 \\
cup & 512 \\
car & 411 \\
bicycle & 171 \\
pizza & 150 \\
apple & 105 \\
dog & 67 \\
cat & 54 \\
\bottomrule
\end{tabular}
\end{table}

\subsection{Models, Thresholds, and Runtime}

Grounding uses Grounding DINO Swin-B [2] with box threshold $=0.30$ and text threshold $=0.25$. Segmentation uses SAM 2.1 Hiera-L [4]. Matching uses IoU threshold 0.10.

Execution was performed in a containerized GPU environment (NVIDIA T4 16 GB VRAM, 8 vCPUs, 56 GB RAM). Randomized steps use seed 42.

\section{Results}

\subsection{Global Prompt Performance}

Simple prompts yield mIoU 0.527 [0.515--0.539] (micro), 0.506 [0.489--0.522] (macro), and 0.513 [0.477--0.547] under strict class+instance balancing.

These results show that the main conclusion (simple prompts outperforming alternatives) holds in both scenarios: natural distribution (production-realistic) and class-balanced distribution (fairness-oriented).

\begin{table}[h]
\caption{Prompt formulation comparison (global summary)}
\label{tab:prompt-comparison}
\centering
\begin{tabular}{lccc}
\toprule
Prompt & mIoU (micro) & mIoU (macro) & Detection (micro) \\
\midrule
Simple & 0.527 & 0.506 & 66.3\% \\
Contextual & 0.250 & 0.366 & 30.7\% \\
Direct & 0.210 & 0.325 & 25.4\% \\
Object-prefixed & 0.204 & 0.305 & 24.8\% \\
\bottomrule
\end{tabular}
\end{table}

Image-level paired Wilcoxon tests reject equality for all comparisons vs. simple prompts: direct ($p=1.08\times10^{-66}$, $\delta=0.522$), contextual ($p=6.95\times10^{-55}$, $\delta=0.427$), and object-prefixed ($p=5.40\times10^{-65}$, $\delta=0.523$).

\begin{figure}[t]
\centering
\includegraphics[width=0.92\linewidth]{../results/figures/e3_prompt_interaction.png}
\caption{Prompt-type interaction analysis.}
\label{fig:prompt-interaction}
\end{figure}

\subsection{Class-level Prompt Effects}

Table~\ref{tab:per-class} shows per-class mean mask IoU by prompt type. The simple prompt is best across all 10 classes, with the strongest absolute gains in person, bottle, and car.

\begin{table*}[t]
\caption{Per-class mean mask IoU by prompt type (EN-only).}
\label{tab:per-class}
\centering
\scriptsize
\begin{tabular}{lcccc}
\toprule
Class & Simple & Contextual & Direct & Object-prefixed \\
\midrule
apple    & 0.260 & 0.200 & 0.122 & 0.099 \\
bicycle  & 0.303 & 0.262 & 0.233 & 0.217 \\
bottle   & 0.558 & 0.279 & 0.189 & 0.152 \\
car      & 0.476 & 0.225 & 0.186 & 0.179 \\
cat      & 0.826 & 0.808 & 0.805 & 0.727 \\
chair    & 0.337 & 0.200 & 0.183 & 0.187 \\
cup      & 0.372 & 0.294 & 0.196 & 0.156 \\
dog      & 0.767 & 0.749 & 0.724 & 0.714 \\
person   & 0.667 & 0.207 & 0.186 & 0.198 \\
pizza    & 0.490 & 0.437 & 0.427 & 0.422 \\
\bottomrule
\end{tabular}
\end{table*}

\begin{figure}[t]
\centering
\includegraphics[width=0.92\linewidth]{../results/figures/e1_class_comparison.png}
\caption{Per-class comparison across prompt formulations.}
\label{fig:class-comparison}
\end{figure}

\subsection{Grounding vs. Segmentation Bottlenecks}

Failure decomposition shows 73.8\% grounding misses and only 2.17\% pure segmentation errors among matched detections. When localization is strong, segmentation remains robust (conditional mask IoU 0.849 at box IoU $\geq$ 0.75).

Failure taxonomy is: 73.8\% grounding misses, 16.2\% false positives, 7.4\% incomplete masks, and 1.8\% incomplete boxes.

\begin{figure}[t]
\centering
\includegraphics[width=0.92\linewidth]{../results/figures/e5_failure_taxonomy.png}
\caption{Failure taxonomy across all evaluated instances.}
\label{fig:failure-taxonomy}
\end{figure}

\begin{figure}[t]
\centering
\includegraphics[width=0.48\linewidth]{../results/figures/e5_failures/failure_00.png}
\hfill
\includegraphics[width=0.48\linewidth]{../results/figures/e5_failures/failure_01.png}
\caption{Qualitative failure examples illustrating grounding misses and low-quality mask outcomes.}
\label{fig:qualitative-failures}
\end{figure}

\section{Discussion}

Results support a consistent \emph{less is more} effect for prompt formulation in this modular setup. Three mechanisms remain plausible: (i) dilution of class token salience by extra functional tokens, (ii) train-inference mismatch between grounding pretraining and instruction-like prompts, and (iii) semantic mismatch between segmentation-intent wording and grounding objectives.

From a machine-vision systems perspective, practical optimization should prioritize the grounding stage (prompt policy, class-wise thresholding, grounding backbone upgrades) before replacing the segmentation module.

From an evaluation standpoint, class standardization does not require re-running full model inference for the current dataset: balancing is applied at analysis level using the already generated instance-level outputs. A full re-execution becomes methodologically necessary only if a higher class floor is desired (e.g., 100 or 150 instances per class), which requires a new sample with additional images because the current minimum class count is 54.

Limitations include a single benchmark (COCO val2017), fixed thresholds, EN-only primary analysis, and absence of direct human-time annotation studies. Future work should include threshold sweeps, multilingual robustness, and cross-dataset validation.

\section{Threats to Validity}

\textbf{Internal validity.} The main risk is metric sensitivity to matching protocol and thresholds. We mitigate this by using fixed matching rules, one-to-one Hungarian assignment, and reporting both aggregate and image-level inferential results.

\textbf{Construct validity.} Pre-labeling utility is approximated through overlap-based metrics and utility buckets, not through direct annotation-time measurement. This may under- or over-estimate practical productivity gains in human workflows.

\textbf{External validity.} Results are reported on COCO val2017 and one modular model pair (Grounding DINO + SAM 2.1). Generalization to domain-shifted datasets, multilingual prompts, and alternative grounding/segmentation combinations should be validated in future experiments.

\section{Future Work}

Three directions are prioritized for follow-up studies.

\textbf{Larger class-balanced protocols.} The current class-balanced analysis is capped by the smallest class support (54 instances/class). A next step is to re-sample and re-run inference with higher class floors (e.g., 100--150 instances/class) to increase statistical power for minority classes.

\textbf{Model and threshold sensitivity.} We plan controlled sweeps over grounding thresholds and alternative grounding backbones while keeping the segmentation stage fixed, followed by ablations where SAM 2.1 variants are swapped under identical grounding outputs.

\textbf{Operational validation with human annotators.} Beyond overlap metrics, we plan a human-in-the-loop study measuring annotation time, correction effort, and inter-annotator consistency across prompt formulations to quantify real productivity gains.

\section{Conclusion}

Prompt formulation is a high-impact lever in modular vision-language pre-labeling pipelines. In this study, single class-name prompts outperform more elaborate alternatives across micro, macro, balanced, and image-level inferential views. Performance degradation is predominantly due to grounding failures, while segmentation quality remains high once localization is correct.

The main operational recommendation is clear: default to simple prompts, optimize grounding first, and adopt class-level ROI triage for deployment planning.

\section*{Declarations}

\textbf{Funding} No funding was received for conducting this study.

\textbf{Competing Interests} The authors declare no competing interests.

\textbf{Author Contributions} Conceptualization, methodology, implementation, experiments, data curation, formal analysis, and manuscript drafting were performed by the author team and reviewed jointly. All authors approved the final manuscript.

\textbf{Data Availability} Processed experimental artifacts, prompt templates, and analysis tables are available at \url{\publicrepo}. A release archive DOI will be added upon public release finalization.

\textbf{Code Availability} Pipeline source code and statistical analysis scripts are available at \url{\publicrepo}. Reproducible release tags will be provided in the public repository.

\textbf{Ethics Approval} Not applicable.

\textbf{Consent to Participate} Not applicable.

\textbf{Consent for Publication} Not applicable.

\begin{thebibliography}{99}

\bibitem{ref1} Lin TY, Maire M, Belongie S, et al. Microsoft COCO: Common Objects in Context. In: \emph{European Conference on Computer Vision (ECCV)}. 2014.
\bibitem{ref2} Liu S, Zeng Z, Ren T, et al. Grounding DINO: Marrying DINO with Grounded Pre-Training for Open-Set Object Detection. arXiv:2303.05499. 2023.
\bibitem{ref3} Kirillov A, Mintun E, Ravi N, et al. Segment Anything. In: \emph{IEEE/CVF International Conference on Computer Vision (ICCV)}. 2023.
\bibitem{ref4} Ravi N, Gabeur V, Hu YT, et al. SAM 2: Segment Anything in Images and Videos. arXiv:2408.00714. 2024.
\bibitem{ref5} Ganguly A, et al. Labeling Copilot: A Deep Research Agent for Automated Data Curation in Computer Vision. arXiv:2509.22631. 2025.
\bibitem{ref6} Comparison of Manual and AI-Assisted Labeling Techniques in Pixel-Wise Instance Segmentation. IEEE Xplore Document 10883307.
\bibitem{ref7} Zhou X, et al. Open-Vocabulary Object Detection in UAV Imagery: A Review and Future Perspectives. 2025.
\bibitem{ref8} Liang F, Wu B, Dai X, et al. Open-Vocabulary Semantic Segmentation with Mask-Adapted CLIP. In: \emph{IEEE/CVF Conference on Computer Vision and Pattern Recognition (CVPR)}. 2023.
\bibitem{ref9} Acuna D, Kar A, Fidler S. Efficient Interactive Annotation of Segmentation Datasets with Polygon-RNN++. In: \emph{IEEE/CVF Conference on Computer Vision and Pattern Recognition (CVPR)}. 2018.
\bibitem{ref10} Radford A, Kim JW, Hallacy C, et al. Learning Transferable Visual Models From Natural Language Supervision. In: \emph{International Conference on Machine Learning (ICML)}. 2021.
\bibitem{ref11} Zhou K, Yang J, Loy CC, Liu Z. Learning to Prompt for Vision-Language Models (CoOp). \emph{International Journal of Computer Vision}. 2022.
\bibitem{ref12} Lu Y, Liu X, Zhang T, et al. Prompt Distribution Learning (ProDA) for Vision-Language Models. arXiv:2205.03340. 2022.
\bibitem{ref13} Zhong Y, Yang J, Zhang P, et al. RegionCLIP: Region-Based Language-Image Pretraining. In: \emph{IEEE/CVF Conference on Computer Vision and Pattern Recognition (CVPR)}. 2022.
\bibitem{ref14} Zhu X, et al. APE: Aligning Prompt Engineering for Open-Set Detection and Grounding. arXiv preprint. 2023.
\bibitem{ref15} Kuhn HW. The Hungarian Method for the Assignment Problem. \emph{Naval Research Logistics Quarterly}. 1955.
\bibitem{ref16} Munkres J. Algorithms for the Assignment and Transportation Problems. \emph{Journal of the Society for Industrial and Applied Mathematics}. 1957.
\bibitem{ref17} Carion N, Massa F, Synnaeve G, et al. End-to-End Object Detection with Transformers (DETR). In: \emph{European Conference on Computer Vision (ECCV)}. 2020.

\end{thebibliography}

\end{document}
